\title{LongRoPE: Extending LLM Context Window Beyond 2 Million Tokens}

\begin{abstract}
Expanding the context window in large language models (LLMs) is highly sought after. Nevertheless, this advancement faces obstacles such as substantial fine-tuning expenses, the limited availability of sufficiently lengthy textual data, and the detrimental impact of newly introduced token positions with extreme values. As a result, the current maximum context window for extended LLMs remains capped at approximately 128k tokens.

This work presents LongRoPE, a novel method that, for the first time, expands the context length of pre-trained LLMs to a remarkable \textbf{2048k} tokens. This is accomplished with as few as 1k fine-tuning steps at a training length capped at 256k, all while preserving the performance within the original short context window. The approach relies on three primary innovations: (i) we discover and leverage two types of positional interpolation non-uniformities through an efficient search process, which yields improved fine-tuning initialization and supports up to an 8$\times$ extension even without fine-tuning; (ii) we propose a two-stage progressive extension technique, whereby the model is first fine-tuned on sequences of length 256k, followed by an additional round of positional interpolation on the already-extended LLM to ultimately reach a 2048k context size; (iii) we perform an 8k-length adjustment of LongRoPE to fully restore performance on short context window tasks. Comprehensive experiments involving LLaMA2 and Mistral over a range of benchmarks attest to the robustness and efficacy of our strategy. The LongRoPE approach requires only minimal changes to the positional embedding component, retaining the rest of the model’s architecture and compatibility with nearly all existing optimizations. Code will be released at \texttt{https://github.com/microsoft/LongRoPE}

\section{Introduction}

Although Large Language Models (LLMs) have demonstrated impressive capabilities across a range of tasks~\cite{gpt4,llama2}, they are typically constrained by a finite context window, such as the 4096-token cap in LLaMA2~\cite{llama2}. Once the input surpasses this limit, LLMs exhibit diminished performance, as they encounter positional information not covered during training. This constraint creates difficulties for several key applications, including in-context learning with a substantial number of examples~\cite{Huang2023FewerIM}, as well as for LLM-driven agent systems~\cite{park2023generative,madaan2023selfrefine}.

Recent research demonstrates that the context window of pre-trained LLMs can be expanded to approximately 128k tokens through fine-tuning on larger text corpora~\cite{longlora,pi,yarn,zhang2024soaring,liu2023scaling}. However, the process faces three primary challenges when attempting even greater context expansion. Firstly, the introduction of numerous unfamiliar position indices—those which were not trained during the original pre-training—results in an abundance of anomalous values, which generates out-of-distribution complications and hampers the convergence of fine-tuning~\cite{pi}. This issue is pronounced when the context length jumps dramatically, as moving from 4k to over 1000k brings in over 90\% novel positions. Secondly, fine-tuning typically necessitates input texts with lengths matching the desired context window, but existing datasets seldom contain texts exceeding 1000k tokens. Additionally, training on such exceptionally lengthy sequences demands substantial computational power, leading to excessive time and GPU consumption. Thirdly, scaling context windows to extreme lengths weakens attention mechanisms since attention is distributed across many token positions, thereby diminishing performance for tasks within shorter, original contexts~\cite{pi}.

A common strategy to address the initial challenge involves interpolating RoPE positional embeddings~\cite{rope,pi}, where novel position indices are scaled down to fit within the original pretrained boundaries, as depicted in Fig.\ref{fig:overview}. Position Interpolation (PI)~\cite{pi} performs a linear interpolation of RoPE's rotational angles based on the extension factor. NTK~\cite{ntk,dynamicntk}, by contrast, proposes uneven interpolation and extrapolation methods for different RoPE dimensions. YaRN~\cite{yarn} further divides RoPE dimensions according to frequency and applies three distinct techniques: extrapolation, NTK interpolation, and linear interpolation, each to its designated group. Nevertheless, the positional embedding within Transformer models displays a \emph{complex, non-uniform distribution of information entropy}. Current methodologies fail to harness this intricate non-uniformity, which results in information degradation and ultimately constrains the achievable context window length.

Section~\ref{sec:analysis} presents two major empirical observations: \textbf{(1)} Robust positional interpolation must address two types of non-uniformities, namely, the varying RoPE dimensions and the distribution of token positions. Less interpolation is beneficial for both lower RoPE dimensions and for tokens occurring earlier in the sequence; nonetheless, the most suitable strategies vary with the intended extension length. \textbf{(2)} Incorporating these sources of non-uniformity into the interpolation process aids in preserving crucial information from the original RoPE—particularly for essential dimensions and specific token positions. This approach reduces the adverse effects typically introduced by interpolation and supplies an improved initialization for subsequent fine-tuning. Additionally, it supports up to an 8$\times$ extension when fine-tuning is not applied.

Building on these results, we present \textbf{LongRoPE}, a robust approach designed to increase the context window in LLMs to more than 2 \emph{million} tokens. The foundation of LongRoPE lies in three principal innovations. To start, {LongRoPE} capitalizes on multidimensional positional interpolation, leveraging non-uniformities in these dimensions. Specifically, it determines optimal rescale factors for RoPE's rotation angles in each dimension according to token locations. Since the space to discover appropriate rescale factors grows exponentially relative to the desired window expansion, {LongRoPE} employs an evolutionary search algorithm supplemented with two distinct optimization methods to maximize the efficiency of the search. Fig.~\ref{fig:overview} provides a representative instance of the rescaled RoPE identified via this search process.

Subsequently, {LongRoPE} employs an efficient, incremental extension methodology to obtain a 2048k context window, circumventing the necessity for explicit fine-tuning on extremely lengthy texts, which are both rare and difficult to access. The process initiates by identifying a 256k length that the pre-trained LLM can handle, followed by fine-tuning at this specific length. Given that our non-uniform positional interpolation technique supports an 8$\times$ expansion even without fine-tuning, we proceed to perform an additional search for updated RoPE rescale factors on the LLM that has already been fine-tuned and extended. As a result, this approach enables a 2048k context window in both LLaMA2 and Mistral~\cite{mistral}.

In order to reduce the decline in performance within the initial (smaller) context window, {LongRoPE} refines the RoPE rescale parameters for the enlarged LLM. Analogous to enlarging the context window from 256k to 2048k, we employ our search method to lower the context windows to 4k and 8k for the 256k fine-tuned LLM, thereby facilitating minimal positional interpolation. When running inference, if the sequence length is below 8k, RoPE is modified with the rescale factors identified through our search.

Comprehensive experimentation with multiple LLMs and a range of long-context tasks confirms the efficacy of our approach. Our findings indicate that LongRoPE consistently sustains low perplexity for evaluation lengths spanning from 4k to 2048k, surpasses 90\% accuracy in passkey retrieval, and performs competitively on established benchmarks within a 4096-token context window. Moreover, LongRoPE is universally compatible with all LLMs utilizing RoPE embedding. We intend to make both our code and LongRoPE-2048k models publicly available.

\section{Non-uniformity in Positional Interpolation}
\label{sec:analysis}

\subsection{Preliminary}

Transformer models require explicit positional information, often in the form of position embedding, to represent the order of input tokens. Our work focuses on the RoPE~\cite{rope} position embedding, which is widely used in recent  LLMs. For a token at position index $n$, its corresponding RoPE encoding can be simplified as follows:

\begin{equation}
		\fontsize{7.8}{7.8} \selectfont
	\label{eq:rope}
	\left[ cos(n\theta_0), sin(n\theta_0), cos(n\theta_1), \cdot\cdot\cdot, cos(n\theta_{d/2-1}), sin(n\theta_{d/2-1}) \right]   
\end{equation}

Here, $d$ denotes the embedding size, $n\theta_i$ specifies the rotational phase applied to the token located at position $n$, and $\theta_i=\theta^{-2i/d}$ corresponds to the frequency for rotation. Typically, RoPE utilizes a base value of $\theta$ set to 10000.

\textbf{\emph{Context window extension ratio $s$} and positional interpolation}. We introduce $s$ as the quotient between the increased context window length $L'$ and the initial length $L$, formulated as $s=\frac{L'}{L}$.

To extend the context window from $L$ to $L'$, current positional interpolation methods suggest downscaling rotation frequencies $\theta_i$ based on extension ratio $s$. Let $\beta=\theta^{2/d}$, and $\lambda$ denote the actual rescale factor related to $s$, we   unify these positional interpolation methods as follows:

\begin{equation}	
	\fontsize{7.5}{7.5} \selectfont
	\label{eq:unified_rope}
	\begin{aligned}
		\left[cos\left(\frac{n}{\lambda(\beta)^0}\right), sin \left(\frac{n}{\lambda(\beta)^0}\right), cos\left(\frac{n}{\lambda(\beta)^1}\right), \cdot\cdot\cdot,  sin \left(\frac{n}{\lambda(\beta)^{d/2-1}}\right) \right]   
	\end{aligned}
\end{equation}

\textbf{Linear positional interpolation (PI)}. PI~\cite{pi} proposes to interpolate position indices in a linear manner, constrained to the maximum sequence length used during pre-training. When expanding to a new sequence length with an extension ratio $s$, the model scales down the rotational angles for every position in the RoPE encoding by a factor of $\lambda=s$. Yet, this compression of positional information leads to a denser representation, making it harder for the model to differentiate between tokens at adjacent positions. As a consequence, PI's effectiveness diminishes as the extension ratio increases.

\textbf{NTK-based interpolation and extrapolation}. ~\cite{ntk,dynamicntk} approach RoPE from the lens of information representation and utilize the Neural Tangent Kernel (NTK) framework~\cite{jacot2018neural,tancik2020fourier}. To address the problem of closely clustered positional indices in Positional Interpolation (PI), these works propose to distribute the interpolation burden more evenly across the various dimensions of RoPE. By applying smaller scaling factors to the lower (i.e., higher frequency) dimensions and amplifying the scaling in higher (i.e., lower frequency) dimensions, their method facilitates both the interpolation and extrapolation of positions, where $\lambda=s^{i}$. The refined dynamic NTK model~\cite{dynamicntk} further tunes the extension ratio at each sequence position in accordance with the current sequence's length. In contrast to PI, which requires model fine-tuning, NTK-aware approaches allow for the extension of context windows without additional training; however, they typically support up to a $4\times$ extension ratio.

\textbf{YaRN}~\cite{yarn} substantially enhances the efficacy of positional interpolation by organizing RoPE dimensions into three distinct groups, classified by their frequencies. For dimensions with high frequency, an extrapolation approach is utilized ($\lambda$=1). Linear interpolation (PI) is applied to those with low frequencies, whereas the intermediate frequency dimensions adopt the NTK method. The core innovation of YaRN is the method it uses to cluster RoPE dimensions; however, this process presently relies on manual, empirical selection. As a result, it may not yield optimal outcomes when applied to newly developed LLMs.

\subsection{Study on Non-uniform Positional Interpolation}

Drawing upon the ideas from NTK and YaRN, we observe their improvement stems from leveraging non-linear properties, particularly by accounting for varying frequencies in RoPE dimensions for targeted interpolation and extrapolation. Nevertheless, existing approaches to non-linearity remain largely constrained by manually crafted heuristics. This observation prompts two fundamental inquiries: (1) Does the present method for positional interpolation represent the best possible approach? (2) Are there additional forms of non-linearity yet to be discovered?

In order to address these issues, we employ evolution search (refer to Sec.~\ref{sec:method}) to identify more effective non-uniform positional interpolations tailored for LLaMA2-7B. The process is steered by perplexity metrics computed over 5 randomly chosen samples from the PG19~\cite{pg19} validation dataset. Our experimental investigation leads us to several important conclusions.

\textbf{Finding 1}: \textit{RoPE dimensions exhibit substantial non-uniformities, which are not effectively handled by current positional interpolation methods.}

We determine the optimal $\lambda$ parameter for each RoPE dimension as specified in Eq.~\ref{eq:unified_rope}. Table~\ref{tbl:exp1} presents a comparison of LLaMA2-7B perplexities achieved by various approaches on the PG19 and Proof-pile~\cite{proof-pile} test datasets, all evaluated without further fine-tuning. Our optimized approach yields notable perplexity reductions, indicating that both the linear (PI) and non-uniform (Dynamic-NTK and YaRN) interpolation techniques do not reach optimality. Of particular interest, YaRN performs worse than PI and NTK on PG19; this may be due to its inability to extend to the intended context window length in language models that have not been fine-tuned. To illustrate, when employing YaRN with an 8k context window, a marked increase in perplexity is observed after 7k tokens.

Our investigation revealed that the rescaled factors $\lambda$ in Eq.~\ref{eq:unified_rope} are non-uniform, unlike the constant scaling parameter $s$ used in PI, NTK’s approach, or the group-wise computation adopted by YaRN. Employing these non-uniform $\lambda$ values yields substantial enhancements in LLaMA2's language modeling capability (as measured by perplexity) over 8k and 16k context lengths, all without necessitating any fine-tuning. The improved results stem from the new positional embeddings' ability to better retain key features of the original RoPE, with a particular emphasis on crucial dimensions, which subsequently decreases the model’s challenge of differentiating between nearby token positions.

\textbf{Finding 2}: \textit{RoPE for the initial tokens in the input sequence should be extrapolated with less interpolation.} 

For the first $\hat{n}$ tokens in the input sequence, we propose that their RoPE should utilize minimal interpolation. The rationale for this stems from their high attention scores, indicating their pivotal role within the attention mechanism, a pattern also noted in Streaming LLM~\cite{streamingllm} and LM-Infinite~\cite{han2023lminfinite}. To assess this idea, we expanded the context window to 8k and 16k using PI and NTK approaches, while preserving the initial  $\hat n$ tokens (where $\hat n$ ranges from 0 to 2 up to 256) without applying interpolation. Notably, when $\hat n = 0$, the configuration reduces to standard PI and NTK. The findings, summarized in Table~\ref{tbl:tokenposition}, point to two main insights: \textbf{(1)} performance benefits from leaving the earliest tokens without interpolation. \textbf{(2)} The ideal quantity of initial tokens $\hat n$ to exclude from interpolation is influenced by the length to which the context window is extended.

\textbf{Finding 3}: \textit{Non-uniform positional interpolation effectively extends LLM context window in both fine-tuning and non-fine-tuning settings.} 

Although our experiments demonstrate that the non-uniform position interpolation discovered through our search yields notable gains in extension capabilities at 8k and 16k without any fine-tuning, extending to even larger contexts necessitates further fine-tuning. Consequently, we perform fine-tuning on LLaMA2-7B using the RoPE configuration identified through our search for a context window of 64k (refer to Appendix for configuration details). As evidenced in Table~\ref{tbl:finetune}, our technique substantially surpasses both PI and YaRN, both prior to and following fine-tuning of LLaMA2-7B. The superiority of our approach arises from its robust exploitation of non-uniform positional interpolation, which reduces information degradation and offers a more effective starting point for fine-tuning.

\textbf{Summary}. Our investigation reveals two important sources of non-uniformity: differentiation across RoPE dimensions and among token positions. Harnessing these non-uniformities through tailored positional interpolation significantly advances LLMs’ performance on extended context windows.

\section{LongRoPE}

\label{sec:method}
Building upon these insights, we propose LongRoPE, a method that initially develops an effective search algorithm to take full advantage of the two forms of non-uniformity. This approach is subsequently employed to expand the LLM context window to support over 2 million tokens.

\subsection{Problem Formulation}

These dual non-uniformities significantly expand the range of possible solutions and add layers of complexity to the optimization process. To manage this challenge, we reformulate the optimization of multidimensional non-uniform position interpolation as a search problem.

For a LLM targeting a  context window size of $L'$ and lengthy input documents $\mathbf{X}$, where each $\mathbf{x}\in\mathbf{X}$ surpasses $L'$ in token length, we denote the original rotary angle of the $i^{th}$ dimension in RoPE embedding at token position $n$ as  $\frac{n}{\beta^i}$. The optimization problem is then formulated as follows:

\begin{equation}
\begin{gathered}
		\label{eq:problem}

\underset{\mathbf{x}\in\mathbf{X}; \, |\mathbf{x}|\ge L'}{ \operatorname{arg\,min} }\, \mathcal{L}\left(\text{LLM}(\text{RoPE}, \mathbf{X})\right), \;\text{with}

\\
\scalebox{0.93}{
$\underset {\begin{subarray}{l} i=0,\ldots,\frac{d}{2}-1;\\ n\in[0,|\mathbf{x}|);\end{subarray}}{\text{RoPE}_ (n)} = {\left[\ldots, \cos\left(\mathbb{I}(\hat{\lambda}_i, \hat{n}) \frac{n}{\beta^i}\right), \sin\left(\mathbb{I}(\hat{\lambda}_i, \hat{n}) \frac{n}{\beta^i}\right), \ldots\right]}$}\\
	\text{where} \; \mathbb{I}(\hat{\lambda}_i, \hat{n}) = \begin{cases}
	1 &\text{if } n < \hat{n}\\
	{\frac{1}{\lambda}_i} &\text{if } n \geq \hat{n}
\end{cases}
	\end{gathered}
\end{equation}

We propose a collection of rescale factors, $\mathbb{I}(\hat{\lambda}_{i},\hat{n})$, to account for two distinct types of non-uniformity. Here, $\hat{\lambda_i}$ represents irregularities across RoPE dimensions, while $\hat{n}$ pertains to non-uniformity in token positions. The function $\mathbb{I}(\hat{\lambda}_{i},\hat{n})$ is employed to adjust the rotation angle in the $i^{th}$ RoPE dimension, with $\hat\lambda_{i}$ serving as the dimension-specific rescale factor and $\hat{n}$ functioning as the position threshold. For the first $\hat{n}$–1 token positions, rescaling via $\hat\lambda_{i}$ is omitted, thus the original RoPE rotary angle $\frac{n}{\beta^i}$ remains unaltered. Conversely, for tokens positioned at $n \ge \hat{n}$, the rescaling factor is utilized.

Suppose a desired context window length $L'$ is specified; the task then becomes selecting the most effective set of rescale factors ($\mathbb{I}(\hat{\lambda}_{0},\hat{n})$, $\mathbb{I}(\hat{\lambda}_{1},\hat{n})$, ..., $\mathbb{I}(\hat{\lambda}_{i},\hat{n})$, ...) for the $1^{st}$ through $d^{th}$ dimensions of RoPE. The aim is for the resulting $\text{LLM}$, utilizing these rescaled $\text{RoPE}$ values, to minimize the next token prediction loss $\mathcal{L}$ (which also corresponds to perplexity) on sequences $\mathbf{X}$ of length $L'$.

\subsection{Searching the Non-uniform Position Interpolation}

To address the challenge presented in Eq.~\ref{eq:problem}, we present a straightforward and robust approach that aims to identify the most suitable RoPE rescale factors, thereby leveraging the diverse non-uniform characteristics inherent in multidimensional position embeddings.

\textbf{Search space}. A comprehensive search space is constructed to capture both forms of non-uniformity. Table~\ref{tbl:searchspace} presents an overview of this search space. In detail, the approach permits the exploration of individualized rescale factors across each dimension within the RoPE embedding. For ease of configuration, the search is conducted over $\lambda_i$ and $\hat{n}$, as opposed to directly searching $\mathbb{I}(\hat{\lambda}_{i},\hat{n})$; here, $\hat{\lambda}_{i}=1/\lambda_i$. As indicated in Table~\ref{tbl:searchspace}, the range for $\lambda_i$ spans from 1.0 (representing straightforward extrapolation) up to $s\times1.25$ (corresponding to broader interpolation relative to PI), incremented in steps of 0.01. In this expression, $s$ denotes the intended context window expansion ratio.

The parameter $\hat{n}$ determines how many leading token positions preserve the original RoPE embedding instead of applying position interpolation. In practice, we select $\hat{n}$ from the set \{0, 1, 2, 4, 8, 12, 16, 20, 24, 28, 32, 64, 128, 256\}. Setting $\hat{n}=0$ means that every token position employs the rescale factors identified by the search.

\textbf{Evolution-based search}. The search space outlined in Table~\ref{tbl:searchspace} encompasses a vast array of positional interpolation configurations, making it difficult to navigate efficiently. For instance, an extension ratio of $s=4\times$ results in $400^{128/2}\times14 = 4\times10^{167}$ possible combinations. As the extension ratio increases, this space grows at an exponential rate. To efficiently manage this complexity, we implement evolution search~\cite{guo2020single} and employ two optimization strategies that substantially accelerate the search process. The complete procedure is presented in Algorithm~\ref{alg:search}.

\textit{Optimized initial population generation}. Rather than filling the initial population of $P$ rescale factors entirely at random, our approach directly includes the three RoPE rescale factors used by PI, NTK, and YaRN as distinct individuals at the outset. The other $P$-3 individuals are produced by randomly mutating these three baseline rescale factors, each mutation occurring with probability $p$.

\textit{Monotonically non-decreasing constraint}. Once the initial population has been generated, we evaluate the LLM perplexity for every member. To do this, we adjust the target LLM with the relevant RoPE rescale factors, and then calculate the perplexity on input $\mathbf{X}$. The highest-ranked individuals are then selected as parents for the subsequent evolution phase. Nonetheless, due to the breadth of the search space, unsophisticated mutation and crossover procedures may traverse unproductive regions, resulting in excess and inefficient perplexity evaluations. This challenge becomes even more pronounced as $L'$ increases, since each perplexity computation incurs significant inference overhead.

To mitigate this issue, we enforce a monotonic non-decreasing condition on the sampled RoPE rescaled factors, such that $\lambda_i \leq \lambda_{i+1}$. Only those RoPE configurations that meet this requirement are utilized within LLMs for evaluating perplexity, which notably reduces the computational burden associated with the search process. More concretely, we stipulate that $\lambda_i$ should grow monotonically as the RoPE dimension increases (i.e., $i$ = 0,...,63). This monotonic relationship across dimensions is motivated by NTK theory~\cite{jacot2018neural,tancik2020fourier,ntk}, which posits that lower dimensions—corresponding to higher frequencies—necessitate less interpolation (i.e., smaller $\lambda_i$), while higher dimensions—associated with lower frequencies—can accommodate greater interpolation (i.e., larger $\lambda_i$).

\begin{algorithm}[t]
	\small
	\caption{The search algorithm}
	\textbf{Input:} target LLM, input samples $\mathbf{X}$, population size $P$, mutation size $N_1$, crossover size $N_2$, maximum iteration count $\mathcal{T}$, mutation probability $p$\\

	\begin{algorithmic}[1]
		\label{alg:search}
		\STATE Top-k $\leftarrow$ $\phi$;	
		\STATE $\text{P}_0$ $\leftarrow$ \textit{Initialize\_population\_with\_optimization} ($P$, $\mathbf{X}$, $p$); 
		\FOR{$i=1$ \textbf{to} $\mathcal{T}$}
		\STATE \textit{Compute\_perplexity} (LLM, $\text{P}_{i-1}$, $\mathbf{X}$);
		\STATE Top-k $\leftarrow$ \textit{Update\_Topk} (Top-k, $\text{P}_{i-1}$);
		\STATE $\text{P}_{mutation}$ $\leftarrow$ \textit{Mutation\_with\_mono\_constraint} (Top-k, $N_1$, $p$); 
		\STATE $\text{P}_{crossover}$ $\leftarrow$ \textit{Crossover\_with\_mono\_constraint} (Top-k, $N_2$);
		\STATE $\text{P}_i$ $\leftarrow$ $\text{P}_{mutation} \cup \text{P}_{crossover} \cup$ Top-k;
		\ENDFOR
		\STATE Output the individual from Top-k exhibiting the minimal perplexity;
	\end{algorithmic}
\end{algorithm}

\textbf{8$\times$ extension achieved without fine-tuning}. Utilizing evolutionary search, our approach efficiently determines optimal non-uniform RoPE rescale factors, safeguarding crucial dimensions and positions to reduce information degradation due to interpolation. As illustrated in Fig.\ref{fig:non-ft}, this technique permits LLaMA2's context window to be expanded from 4k to 32k absent any fine-tuning procedures. Conversely, prevailing methods like PI, non-uniform NTK, and YaRN result in a considerable increase in perplexity after a 2$\times$ context window extension.

\subsection{Extending LLM Context Window to 2048K}
\label{sec:extendto2048k}

\textbf{Progressive extension to 2048k}. In this section, we present our approach for scaling the context window of pre-trained LLMs from the standard 4k tokens to beyond 2048k tokens. Through the use of our non-uniform positional interpolation technique, it is possible to achieve up to an 8$\times$ context length increase without the need for fine-tuning. However, when aiming for substantially larger expansions (such as 512$\times$), fine-tuning becomes indispensable. One potential strategy involves identifying suitable RoPE rescaling factors for the targeted 2048k context window followed by fine-tuning. Nonetheless, this strategy is hindered by the extremely high computational and training costs involved. Additionally, our observations indicate that performing effective fine-tuning at such extensive scaling ratios remains difficult (refer to Appendix for details).

Fortunately, LongRoPE proves to be robust for both baseline and further fine-tuned extended LLMs. In response, we propose a streamlined, incremental approach that reaches the desired 2048k context window using only 1k fine-tuning steps, with a training length capped at 256k.

$\Diamond$ \textit{Expanding pre-trained LLMs to 256k with LongRoPE search}. Using LLaMA2 as a case study, we perform a search aiming for context window sizes of 128k and 256k. At this phase, the context length is extended by factors of 32$\times$ and 64$\times$, respectively.

$\Diamond$ \textit{Fine-tuning to 256k}. In the subsequent phase, we adapt the pre-trained LLM to handle a 256k context window. Initially, we fine-tune LLaMA2 for 400 steps while applying RoPE rescaled factors corresponding to 128k. Afterward, we update the RoPE rescaled factors to target 256k on the resultant checkpoint and proceed with 600 more fine-tuning steps. This staged procedure is found to be more effective than attempting direct fine-tuning to 256k.

$\Diamond$ \textit{Adapting a previously fine-tuned extended LLM to 2048k with LongRoPE search}. In the last stage, we carry out a subsequent search step utilizing the LLM that has already been fine-tuned for a context length of 256k. Remarkably, this process enables scaling to an ultra-long context window of 2048k, achieved without requiring any additional fine-tuning. Overall, this procedure accomplishes an expansion ratio of $512\times$.

\textbf{Addressing degraded performance in shorter context windows}. Upon expanding the context window to an ultra-long 2048k, a decline in effectiveness emerges within the initial context window. This phenomenon has been documented in prior work on positional interpolation~\cite{pi}, where embedding positions into higher dimensions compresses their distribution in the original window, thereby diminishing the language model's capability. At a 512$\times$ extension rate, position embeddings inside the original 4k context window are especially densely packed.

To address this issue, an additional evolutionary search is conducted on the expanded LLM to fine-tune RoPE rescale parameters specifically for shorter context lengths, such as 4k and 8k. The search range for $\lambda$ is decreased since shorter contexts demand less positional interpolation. Throughout inference, the LLM adapts the RoPE rescale factors accordingly based on the input context length.

\section{LongRoPE}

\label{sec:method}
In response to these insights, we propose LongRoPE—a method that initially implements an effective search algorithm designed to leverage both forms of non-uniformity, and subsequently applies this approach to expand the LLM context window to over 2 million tokens.

\subsection{Problem Formulation}

These two forms of non-uniformity expand the solution space significantly, thereby complicating the optimization process. In response, we reformulate the multidimensional non-uniform position interpolation optimization challenge into a search problem.

For a LLM targeting a  context window size of $L'$ and lengthy input documents $\mathbf{X}$, where each $\mathbf{x}\in\mathbf{X}$ surpasses $L'$ in token length, we denote the original rotary angle of the $i^{th}$ dimension in RoPE embedding at token position $n$ as  $\frac{n}{\beta^i}$. The optimization problem is then formulated as follows:

\begin{equation}
\begin{gathered}
		\label{eq:problem}

\underset {\mathbf{x}\in\mathbf{X}; \, |\mathbf{x}|\ge L'}{ \operatorname {arg\,min} } \,  \mathcal{L}\, \left( \text{LLM} (\text{RoPE}, \mathbf{X})\right), \text{with}
\\
\scalebox{0.93}{
$\underset {\begin{subarray}{l} i=0,\cdot\cdot,\frac{d}{2}-1;\\ n\in[ 0,|\mathbf{x}|);\end{subarray}}{\text{RoPE}_(n)}= {\left[\ldots, \cos\left(\mathbb{I}(\hat{\lambda}_{i}, \hat{n}) \cdot \frac{n}{\beta^i}\right), \sin\left(\mathbb{I}(\hat{\lambda}_{i}, \hat{n}) \cdot \frac{n}{\beta^i}\right), \ldots\right] }$}\\
	\text{where} \quad \mathbb{I}(\hat{\lambda}_{i},\hat{n})=\begin{cases}
	1&\mbox{\quad $n<\hat{n}$}\\
{\frac{1}{\lambda}_{i}}&\mbox{\quad $n\geq\hat{n}$}
\end{cases}
	\end{gathered}
\end{equation}

In this approach, we define a collection of rescale factors, $\mathbb{I}(\hat{\lambda}_{i},\hat{n})$, to account for the two types of non-uniformities present in the system. The parameters $\hat{\lambda_i}$ and $\hat{n}$ represent the variability in the RoPE dimensions and the token position boundaries, respectively. To adjust the rotation angle for the $i^{th}$ RoPE dimension, we incorporate $\mathbb{I}(\hat{\lambda}_{i},\hat{n})$, where $\hat\lambda_{i}$ serves as the specific rescale coefficient and $\hat{n}$ establishes the token position cutoff. For the first $\hat{n}-1$ tokens, the rescale factor $\hat\lambda_{i}$ is not utilized; instead, the default RoPE rotary angle $\frac{n}{\beta^i}$ remains in use. When token positions meet or exceed $n \ge \hat{n}$, the rescale factor is then employed.

For a specified target context window length $L'$, the task is to identify the best set of rescaling factors, $\mathbb{I}(\hat{\lambda}_{0},\hat{n}), \mathbb{I}(\hat{\lambda}_{1},\hat{n}), \ldots, \mathbb{I}(\hat{\lambda}_{i},\hat{n}), \ldots$, corresponding to each of the $1^{st}$ through $d^{th}$ dimensions of RoPE. By adjusting RoPE with these factors in the target $\text{LLM}$, the model aims to minimize the next-token prediction loss, $\mathcal{L}$ (also known as perplexity), on input samples $\mathbf{X}$ whose sequence length is $L'$.

\subsection{Searching the Non-uniform Position Interpolation}

To address the challenge described in Eq.~\ref{eq:problem}, we present a straightforward but powerful approach that identifies the best RoPE rescale factors, enabling comprehensive utilization of the multidimensional variations present in position embeddings.

\textbf{Search space}. Our search space is constructed to encompass both types of non-uniformity, resulting in a broader and more flexible set of candidate configurations. Table~\ref{tbl:searchspace} presents a detailed overview of the search space. Rather than searching over $\mathbb{I}(\hat{\lambda}_{i},\hat{n})$ directly, we simplify the process by targeting the parameters $\lambda_i$ and $\hat{n}$, where $\hat{\lambda}_{i}=1/\lambda_i$, allowing for individualized rescaling of each dimension in the RoPE embedding. As indicated in Table~\ref{tbl:searchspace}, $\lambda_i$ is tunable, with permissible values ranging from 1.0 (which corresponds to simple extrapolation) up to $s\times1.25$ (yielding interpolation beyond that of PI), incremented by steps of 0.01. Here, $s$ represents the ratio by which the context window is to be extended.

The parameter $\hat{n}$ determines how many of the earliest token positions preserve their original RoPE embeddings, foregoing position interpolation. In practice, $\hat{n}$ is selected from the set \{0, 1, 2, 4, 8, 12, 16, 20, 24, 28, 32, 64, 128, 256\}. Setting $\hat{n}=0$ implies that rescaled factors, as determined through the search process, are applied to every token position.

\textbf{Evolution-based search}. The search space defined in Table~\ref{tbl:searchspace} encompasses a wide array of positional interpolation options, making effective navigation a complex task. For instance, when applying a $s=4\times$ extension, the number of possible configurations increases dramatically to $400^{128/2}\times14$ = $4\times10^{167}$. As the extension ratio becomes larger, the search space expands at an exponential rate. To tackle this issue, we implement evolution search~\cite{guo2020single} along with two optimization strategies that significantly enhance the efficiency of the search process. The overall methodology is outlined in Algorithm~\ref{alg:search}.

\textit{Optimized initial population generation}. Rather than relying solely on random initialization for a population of $P$ rescale factors, we explicitly include the RoPE rescale factors associated with PI, NTK, and YaRN as distinct members of the starting population. The other $P-3$ individuals are formed by randomly applying mutations to these three rescale factors, each with probability $p$.

\textit{Monotonically non-decreasing constraint}. Upon creating the initial set of individuals, we evaluate the LLM perplexity for each candidate. This is accomplished by applying the relevant RoPE rescale factors within the target LLM and assessing the perplexity on the input $\mathbf{X}$. The individuals with the top-k lowest perplexity scores are chosen as parents for the evolutionary process. Nonetheless, the expansive nature of the search space means that standard mutation and crossover may lead to suboptimal candidates, thus incurring superfluous perplexity evaluations. This inefficiency is exacerbated when $L'$ is large, due to the substantial computational resources required for each perplexity assessment.

To tackle this challenge, we enforce a non-decreasing monotonicity condition on the sampled RoPE rescaling coefficients: $\lambda_i\le \lambda_{i+1}$. Only those RoPE configurations meeting this criterion are utilized within the LLM for perplexity assessments, greatly curtailing the associated search overhead. In detail, this means $\lambda_i$ must grow monotonically as a function of the RoPE dimension index ($i=0, \ldots, 63$). The rationale for this monotonic structure draws from NTK theory~\cite{jacot2018neural,tancik2020fourier,ntk}, which posits that lower-indexed dimensions—corresponding to higher frequencies—necessitate less interpolation (i.e., smaller $\lambda_i$), whereas higher-indexed dimensions—associated with lower frequencies—permit greater interpolation (i.e., larger $\lambda_i$).

\begin{algorithm}[t]
	\small
	\caption{The search algorithm}
	\textbf{Input:} target LLM, input samples $\mathbf{X}$, population size $P$, mutation size $N_1$, crossover size $N_2$, maximum iterations $\mathcal{T}$, mutation probability $p$\\

	\begin{algorithmic}[1]
		\label{alg:search}
		\STATE Top-k $\leftarrow$ $\phi$;	
		\STATE $\text{P}_0$ $\leftarrow$ \textit{Initialize\_population\_with\_optimization}($P$, $\mathbf{X}$, $p$); 
		\FOR{$i=1$ to $\mathcal{T}$}
		\STATE \textit{Compute\_perplexity}(LLM, $\text{P}_{i-1}$, $\mathbf{X}$);
		\STATE Top-k $\leftarrow$ \textit{Update\_Topk}(Top-k, $\text{P}_{i-1}$);
		\STATE $\text{P}_{mutation}$ $\leftarrow$ \textit{Mutation\_with\_mono\_constraint}(Top-k, $N_1$, $p$);
		\STATE $\text{P}_{crossover}$ $\leftarrow$ \textit{Crossover\_with\_mono\_constraint}(Top-k, $N_2$);
		\STATE $\text{P}_i$ $\leftarrow$ $\text{P}_{mutation}$ $\cup$ $\text{P}_{crossover}$ $\cup$ Top-k;
		\ENDFOR
		\STATE Output the member in Top-k possessing the minimum perplexity;
	\end{algorithmic}
\end{algorithm}

\textbf{8$\times$ context extension without fine-tuning}. By utilizing an evolutionary search, we discover non-uniform RoPE rescale factors that retain crucial dimensions and positions, thereby reducing the risk of information degradation during interpolation. As illustrated in Fig.\ref{fig:non-ft}, our approach successfully broadens LLaMA2's context window from 4k up to 32k without the necessity for fine-tuning. In comparison, previously established techniques like PI as well as non-uniform NTK and YaRN exhibit a sharp rise in perplexity beyond a 2$\times$ extension.

\subsection{Extending LLM Context Window to 2048K}
\label{sec:extendto2048k}

\textbf{Progressive extension to 2048k}. In this section, we present our approach for expanding the context window of pre-trained LLMs beyond the conventional 4k up to more than 2048k tokens. Our experiments indicate that by applying non-uniform positional interpolation, we can realize up to an 8$\times$ extension without the need for fine-tuning. When aiming for much larger extensions (such as 512$\times$), however, fine-tuning becomes indispensable. A commonly adopted strategy involves determining appropriate RoPE rescaling parameters suited for the 2048k context length, followed by model fine-tuning. Nonetheless, this method is hindered by the extremely high computational and resource demands of such training. Additionally, as our experience suggests, effective fine-tuning under such significant extension ratios presents considerable difficulties (see Appendix).

Fortunately, LongRoPE demonstrates efficacy with both base and fine-tuned long-context LLMs. To that end, we propose a streamlined, incremental approach that reaches the intended 2048k context window after only 1k fine-tuning iterations, requiring no more than 256k tokens during training.

$\Diamond$ \textit{Scaling up pre-trained LLMs to 256k with LongRoPE search}. Using LLaMA2 for illustration, we explore context window sizes of 128k and 256k tokens during the search process. This step involves increasing the extension ratio to 32$\times$ and 64$\times$, respectively.

$\Diamond$ \textit{Fine-tuning to 256k}. Subsequently, the pre-trained LLM is fine-tuned to extend its context window to 256k. The process starts with a 400-step fine-tuning of LLaMA2 using RoPE rescaled factors set for 128k. After completion, these factors are switched to those corresponding to 256k, and a further 600 fine-tuning steps are performed on the resultant checkpoint. This approach demonstrates greater efficiency compared to commencing the fine-tuning directly with 256k.

$\Diamond$ \textit{Expanding fine-tuned extended LLM to 2048k using LongRoPE search}. In the last stage, we conduct an additional search step on the LLM that was previously fine-tuned for a context length of 256k. This procedure achieves a significantly expanded context window of 2048k, accomplished without any supplementary fine-tuning. The overall extension factor reaches $512\times$.

\textbf{Restoration of performance in reduced context windows}. When the context window is increased to a substantial 2048k length, a degradation in performance is observed within the initial context window size. This phenomenon has been documented with positional interpolation~\cite{pi}, where embedding positions in expanded dimensions are compressed into a restricted space inside the original context window, consequently diminishing the model's efficacy. Specifically, using an extension ratio of 512$\times$ causes the positions within the standard 4k context window to be excessively condensed.

To address this issue, we conduct an additional evolution search on the extended LLM to fine-tune the RoPE rescale factors specifically for shorter context lengths, such as 4k and 8k. For these shorter spans, we constrain the upper bound for the searched $\lambda$, reflecting the diminished need for positional interpolation. At inference time, the LLM adaptively modifies the relevant RoPE rescale factors.

 \section{Experiments}

\subsection{Setup}
\textbf{Evaluation Tasks and models}. LongRoPE is integrated into LLaMA2-7B and Mistral-7B, and their performance is assessed across three dimensions: (1) measuring the perplexity of large language models with expanded context lengths when processing lengthy documents; (2) conducting the Passkey retrieval task, which evaluates the models' capacity to extract a straightforward passkey from a vast collection of unrelated content; and (3) utilizing established LLM benchmarks under a limited context window of 4096 tokens.

\textbf{Fine-tuning}. For LLaMA2, our approach involves a linear learning rate schedule beginning at 2e-5 and employing a global batch size of 32. The model undergoes 400 fine-tuning steps on the Redpajama~\cite{redpajama} dataset, which is partitioned into 128k-token segments, each bounded by BOS and EOS tokens. Upon completing this phase and obtaining the resulting checkpoint, an additional 600 training steps are performed to extend the context window to 256k tokens. Training at the 128k context length utilizes 8 A100 GPUs through a distributed setup~\cite{cube}, while the larger 256k context requires 16 A100 GPUs. For Mistral, we adopt a fixed learning rate of 1e-6 with a global batch size of 64. In accordance with the YaRN~\cite{yarn} methodology, both 128k and 256k models are trained for 400 steps using Together Computer's Long-Data Collections~\cite{mistraldata}, utilizing a sequence length of 16k, and leveraging 4 A100 GPUs.

\textbf{Search}. When the target window size does not exceed 256k, we set $P$=64, $N_1$=$N_2$=16, $p$=0.3, and $\mathcal{T}$=40, choosing the top-32 candidates for mutation and crossover at each iteration. Perplexity values are computed on 5 randomly selected validation examples from the PG19 dataset, ensuring each meets or exceeds the intended context length. For target windows larger than 512k, the sizes for population, mutation, and crossover operations are all reduced by half. In these cases, perplexity is estimated on 3 random validation samples taken from Pile-Books3~\cite{pile}.

\textbf{Baselines}. In order to achieve a 2048k context window, we conducted fine-tuning on models with existing 128k and 256k context lengths. Consequently, this resulted in the creation of LongRoPE-2048k (ft=128k) for LLaMA2 and LongRoPE-2048k (ft=256k) for Mistral. These four models are evaluated alongside prominent context window extension baselines. In particular, we include publicly available LLMs fine-tuned with positional interpolation strategies such as PI, NTK, and YaRN. Among these are Together-32k~\cite{together}, Code LLaMA~\cite{codellama}, LongLoRA-full-FT-100k~\cite{longlora}, YaRN-LLaMA, and YaRN-Mistral~\cite{yarn}.

\subsection{Main Results}

\label{sec:mainresult}
\textbf{Language modeling performance on sequences up to 256k tokens}. To start, we assess our model in comparison to leading extended LLMs, focusing on a maximum evaluation sequence length of 256k tokens. To highlight the breadth of our approach, we report results on two datasets: the Proof-pile~\cite{pg19} and PG19~\cite{pile} test sets. We compute perplexity across a range of context window sizes, using a sliding window method with a width of 256 tokens. For PG19, evaluation is conducted over the complete test split, comprising 100 documents. For Proof-pile, we adopt the methodology of YaRN~\cite{yarn}, selecting at random 10 samples, each with lengths of at least 128k tokens.

Table~\ref{tbl:proof-pile} and Table~\ref{tbl:pg19} present perplexity results for LLaMA2 and Mistral models augmented with various interpolation strategies on the Proof-pile and PG19 datasets, respectively. Two primary findings emerge: \textbf{(1)} our extended models exhibit a consistent reduction in perplexity as the evaluation length increases from 4k to 256k, demonstrating effective utilization of extended context. \textbf{(2)} Remarkably, even when assessed with a 16$\times$ longer context window—a scenario known to make retaining short-context performance difficult—our LongRoPE-2048k models achieve superior results over leading baselines within the 256k context length range.

\textbf{Modeling language over sequences longer than 2000k}. In order to assess performance with exceptionally lengthy texts, we utilize the Books3~\cite{pile} collection. To streamline evaluation, we randomly pick 20 books, all of which have lengths greater than 2048k, applying a sliding window of 256k tokens for processing.

Table~\ref{tbl:books3} demonstrates that LongRoPE effectively expands the context window of LLaMA2-7B and Mistral-7B models up to 2048k, while also maintaining or improving perplexity relative to baseline approaches within shorter ranges (8k-128k). Substantial performance disparities are apparent between the LLaMA2 and Mistral models at the 2048k context length. Specifically, Mistral shows stronger results at shorter context lengths, but its perplexity rises above 7 when evaluated beyond 256k tokens. LLaMA2's behavior follows anticipated trends: perplexity declines as context length increases, aside from minor upticks at 1024k and 2048k. Notably, LongRoPE-2048k applied to LLaMA2 demonstrates improved performance when fine-tuned with a 256k window compared to 128k, which is attributed to a lower secondary extension ratio (8$\times$ versus 16$\times$). Conversely, Mistral yields better outcomes with a fine-tuning context of 128k. This discrepancy arises because Mistral’s fine-tuning at 128k and 256k adheres to YaRN’s protocol, utilizing a 16k training context that impacts its capacity to extend context window length post fine-tuning.

\textbf{Passkey retrieval}. Next, we examine the functional size of the context window in generative tasks. We utilize a synthetic evaluation procedure for passkey retrieval as described by ~\cite{passkey}. Here, the objective is for the model to locate a randomly inserted passkey—a five-digit numeral—concealed within an extensive document. The structure of the prompt utilized is provided in the appendix. For assessment, we conduct 10 repetitions of the passkey retrieval challenge, with the passkey randomly positioned, following a uniform distribution throughout the length of the context used for evaluation.

Fig.~\ref{fig:passkey} presents a comparative analysis of retrieval accuracy against several baseline models. The performance of existing models diminishes sharply, falling to zero once the token count exceeds 128k. In stark contrast, our LongRoPE-LLaMA2-2048k (ft=256k) demonstrates robust retrieval ability—retaining accuracy at or above 90\% throughout the token range from 4k up to 2048k, even under the demanding scenario of extracting a passkey amid millions of tokens. Meanwhile, LongRoPE-Mistral-2048k (ft=128k) achieves a perfect 100\% accuracy up to 1800k tokens; its accuracy decreases to 60\% at the 2048k mark, which is consistent with the results in Table~\ref{tbl:books3}, where perplexity shows a modest rise at 2048k.

\textbf{Evaluation on conventional benchmarks within the initial context window}. We assess LongRoPE-2048k models based on the Hugging Face Open LLM Leaderboard~\cite{huggingfaceboard}, employing both zero-shot and few-shot methodologies. Specifically, we utilize ARC-Challenge~\cite{allenai:arc} with 25-shot, HellaSwag~\cite{zellers2019hellaswag} with 10-shot, MMLU~\cite{mmlu} with 5-shot, and TruthfulQA~\cite{lin2021truthfulqa} with 0-shot configurations.

Table~\ref{tbl:commonsense} demonstrates that our models perform similarly to the baseline established for a shorter context window, and actually surpass the original Mistral by 0.5\% on the TruthfulQA benchmark. While LongRoPE-LLaMA2-2048k, which was fine-tuned using 256k context length, experiences a marginally greater drop in performance, it generally maintains competitive results across various tasks.

\subsection{Ablation Results}

\textbf{Assessing the Impact of Second Positional Interpolation}. As part of our progressive extension framework, we apply our search algorithm to perform a secondary, non-uniform positional interpolation on the fine-tuned, expanded language models. To assess the benefits of this method, we experiment with a LLaMA2-256k model that has undergone fine-tuning. We then expand this model to 512k, 1024k, and 2048k contexts using both PI and YaRN techniques. Results summarized in Table~\ref{tbl:secondsearch} indicate that our non-uniform positional interpolation maintains perplexity at a stable level. Conversely, when models are extended using PI and YaRN, perplexity rises rapidly as the extension ratio increases.

\textbf{Recovery performance at reduced context sizes.} In order to address the degradation in performance encountered at shorter context windows, we reconfigure the RoPE parameters for LongRoPE-2048k utilizing our search methodology. This involves reducing the upper limit of scale factors allowed during the search process, thereby promoting minimized interpolation when operating at 4k and 8k context lengths. Table~\ref{tbl:shortperformance} presents a perplexity analysis of LongRoPE-LLaMA2-2048k on Proof-pile at these shorter lengths, alongside the mean accuracy achieved across LLM benchmarks. The findings indicate a marked enhancement in performance for short context scenarios.

\textbf{Evaluation of the Impact of Two Types of Non-uniformities}. To understand the role each non-uniformity plays in performance, we conducted ablation studies. Specifically, we designed two experimental setups: (i) adapting LLaMA2-7B to shorter sequence lengths of 16k and 32k through various approaches—including PI, searching solely for the optimal RoPE dimension, and considering both sources of non-uniformity; (ii) scaling our previously fine-tuned LLaMA2 model (trained on 256k-length) up to 2048k using an analogous strategy. Perplexity assessments were performed without additional fine-tuning. The results displayed in Table~\ref{tbl:searchdimension} indicate that incorporating non-uniformity in the RoPE dimension yields a marked reduction in perplexity relative to the PI method’s linear interpolation. Adjusting for non-uniformity in token position provides noticeable benefits at the 16k and 32k lengths, though its efficacy diminishes at 2048k, likely because of the model’s handling of exceptionally long sequences. Merely retaining the initial tokens without interpolation fails to yield meaningful improvements; we plan to address this limitation in future research.

\section{Related Works}

Apart from techniques relying on position interpolation, this section also reviews other relevant strategies.

\textbf{Retrieval-based} methods employ external memory systems to store extensive historical contexts and utilize retrieval mechanisms to access pertinent documents during inference~\cite{longllama,wang2023augmenting,borgeaud2022improving}. Such methods generally require substantial architectural changes to the LLMs. In comparison, our proposed method involves minimal alterations, limited to positional embedding adjustments, making it more efficient. Furthermore, our approach accommodates a broader range of tasks involving lengthy contexts, including tasks such as summarizing long documents and facilitating few-shot learning, rather than focusing exclusively on retrieval.

\textbf{Extensions of the context window through attention mechanisms}. In addition to positional embedding interpolation, several studies have explored methods to enlarge the input context for LLMs while keeping the context window length unchanged by modifying the attention computation~\cite{han2023lminfinite, streamingllm,parallelwindow}. Central to these approaches is the application of innovative attention masks that alleviate the problem of excessive attention overhead associated with novel positions. These strategies serve as a complement to positional interpolation techniques.

\textbf{Fine-tuning based approaches} are concerned with improving the adaptation of pre-trained LLMs to extended sequence lengths by adjusting position embeddings. Approaches such as Code LLaMA~\cite{codellama}, LLaMA2 Long~\cite{xiong2023effective}, and ScaledRoPE~\cite{liu2023scaling} increase the RoPE base considerably and conduct fine-tuning at the desired context length. In contrast, our approach is adaptable to multiple target lengths and can support sequence lengths surpassing 2M. In the context of long context fine-tuning (above 128k), which is computationally intensive, methods like LongLoRA~\cite{longlora} and PoSE~\cite{zhu2023pose} have been proposed to reduce resource demands. Our strategy operates independently and can complement these resource-efficient fine-tuning techniques.

\section{Conclusion}

This study introduces LongRoPE, a technique that significantly broadens the context window of LLMs to an extraordinary 2048k tokens, all while preserving performance in the original, smaller window. By leveraging two distinct kinds of non-uniformities present in RoPE positional embeddings—identified via a streamlined evolutionary search—we attain dual advantages: optimal initialization for subsequent fine-tuning and immediate 8$\times$ expansion of context length without fine-tuning. Further, we implement a gradual extension approach that employs LLMs fine-tuned at 256k length to scale up to a 2048k window without necessitating additional fine-tuning. Comprehensive experimental results affirm the efficacy of LongRoPE. We anticipate that LongRoPE-2048k models will facilitate a wide array of long-context use cases and stimulate new directions in research.

\section*{Broader Impacts} The aim of this research is to contribute to the progression of Machine Learning. Although our work may give rise to a variety of societal effects, we do not believe that any particular outcome needs to be explicitly called out at this time.

	\balance

\end{document}